\documentclass[10pt,a4paper]{amsart}
\usepackage{amsmath,amssymb,amsthm,graphicx}
% Retained minimal preamble for the Stacks Project chapter flat.tex (More on Flatness).
% Verbatim excerpt of the chapter preamble: Xy-pic setup (native preamble lines 14-18) and the native
% theorem declarations with the mathematical macros (native preamble lines 154-217). Copied mechanically;
% nothing added, removed or redefined. The wrapper class and the amsmath/amssymb/amsthm/graphicx base
% packages are supplied by the assembler, so the class/externaldocument part of the native preamble is omitted.
% For commutative diagrams we use Xy-pic
\usepackage[all]{xy}

% We use 2cell for 2-commutative diagrams.
\xyoption{2cell}
% Theorem environments.
%
\theoremstyle{plain}
\newtheorem{theorem}[subsection]{Theorem}
\newtheorem{proposition}[subsection]{Proposition}
\newtheorem{lemma}[subsection]{Lemma}

\theoremstyle{definition}
\newtheorem{definition}[subsection]{Definition}
\newtheorem{example}[subsection]{Example}
\newtheorem{exercise}[subsection]{Exercise}
\newtheorem{situation}[subsection]{Situation}

\theoremstyle{remark}
\newtheorem{remark}[subsection]{Remark}
\newtheorem{remarks}[subsection]{Remarks}

\numberwithin{equation}{subsection}

% Macros
%
\def\lim{\mathop{\mathrm{lim}}\nolimits}
\def\colim{\mathop{\mathrm{colim}}\nolimits}
\def\Spec{\mathop{\mathrm{Spec}}}
\def\Hom{\mathop{\mathrm{Hom}}\nolimits}
\def\Ext{\mathop{\mathrm{Ext}}\nolimits}
\def\SheafHom{\mathop{\mathcal{H}\!\mathit{om}}\nolimits}
\def\SheafExt{\mathop{\mathcal{E}\!\mathit{xt}}\nolimits}
\def\Sch{\mathit{Sch}}
\def\Mor{\mathop{\mathrm{Mor}}\nolimits}
\def\Ob{\mathop{\mathrm{Ob}}\nolimits}
\def\Sh{\mathop{\mathit{Sh}}\nolimits}
\def\NL{\mathop{N\!L}\nolimits}
\def\CH{\mathop{\mathrm{CH}}\nolimits}
\def\proetale{{pro\text{-}\acute{e}tale}}
\def\etale{{\acute{e}tale}}
\def\QCoh{\mathit{QCoh}}
\def\Ker{\mathop{\mathrm{Ker}}}
\def\Im{\mathop{\mathrm{Im}}}
\def\Coker{\mathop{\mathrm{Coker}}}
\def\Coim{\mathop{\mathrm{Coim}}}

% Boxtimes
%
\DeclareMathSymbol{\boxtimes}{\mathbin}{AMSa}{"02}

%
% Macros for moduli stacks/spaces
%
\def\QCohstack{\mathcal{QC}\!\mathit{oh}}
\def\Cohstack{\mathcal{C}\!\mathit{oh}}
\def\Spacesstack{\mathcal{S}\!\mathit{paces}}
\def\Quotfunctor{\mathrm{Quot}}
\def\Hilbfunctor{\mathrm{Hilb}}
\def\Curvesstack{\mathcal{C}\!\mathit{urves}}
\def\Polarizedstack{\mathcal{P}\!\mathit{olarized}}
\def\Complexesstack{\mathcal{C}\!\mathit{omplexes}}
% \Pic is the operator that assigns to X its picard group, usage \Pic(X)
% \Picardstack_{X/B} denotes the Picard stack of X over B
% \Picardfunctor_{X/B} denotes the Picard functor of X over B
\def\Pic{\mathop{\mathrm{Pic}}\nolimits}
\def\Picardstack{\mathcal{P}\!\mathit{ic}}
\def\Picardfunctor{\mathrm{Pic}}
\def\Deformationcategory{\mathcal{D}\!\mathit{ef}}

% --- appended 2026-10-12 (Chief): the project own definitions of the reference/slogan/history
% display environments, copied mechanically from the shared Stacks preamble lines 8-12
% (\usepackage{verbatim} plus the three \newenvironment lines, which the chapter preamble keeps
% verbatim from the project preamble). Nothing added, removed or redefined: these environments
% carry no mathematical content and the project itself typesets them as comments. Without these
% three lines a counted statement whose own text uses them fails with Environment slogan undefined.
\usepackage{verbatim}
\newenvironment{reference}{\comment}{\endcomment}
\newenvironment{slogan}{\comment}{\endcomment}
\newenvironment{history}{\comment}{\endcomment}
\usepackage{cleveref}
\crefname{originalreference}{original reference}{original references}
\Crefname{originalreference}{Original reference}{Original references}
\makeatletter
\newlabel{situation-complex-and-divisor}{{0F8S}{source}{}{sourcepointer-1}{}}
\newlabel{situation-complex-and-divisor@cref}{{[originalreference][0][]0F8S}{[1][0][]source}{}{}{}}
\makeatother
\crefname{originallemma}{original lemma}{original lemmas}
\Crefname{originallemma}{Original lemma}{Original lemmas}
\makeatletter
\newlabel{cohomology-lemma-Leta}{{0F8Q}{source}{}{sourcepointer-2}{}}
\newlabel{cohomology-lemma-Leta@cref}{{[originallemma][0][]0F8Q}{[1][0][]source}{}{}{}}
\makeatother
\crefname{originallemma}{original lemma}{original lemmas}
\Crefname{originallemma}{Original lemma}{Original lemmas}
\makeatletter
\newlabel{perfect-lemma-affine-compare-bounded}{{06Z0}{source}{}{sourcepointer-3}{}}
\newlabel{perfect-lemma-affine-compare-bounded@cref}{{[originallemma][0][]06Z0}{[1][0][]source}{}{}{}}
\makeatother
\crefname{originallemma}{original lemma}{original lemmas}
\Crefname{originallemma}{Original lemma}{Original lemmas}
\makeatletter
\newlabel{more-algebra-lemma-ideal-well-defined}{{0GSS}{source}{}{sourcepointer-4}{}}
\newlabel{more-algebra-lemma-ideal-well-defined@cref}{{[originallemma][0][]0GSS}{[1][0][]source}{}{}{}}
\makeatother
\crefname{originallemma}{original lemma}{original lemmas}
\Crefname{originallemma}{Original lemma}{Original lemmas}
\makeatletter
\newlabel{lemma-eta-stalks}{{0F8M}{source}{}{sourcepointer-5}{}}
\newlabel{lemma-eta-stalks@cref}{{[originallemma][0][]0F8M}{[1][0][]source}{}{}{}}
\makeatother
\title{For a good triple a locally bounded complex of finite locally free modules representing the twist exists whenever one exists for the original}
\author{The Stacks Project Authors}
\date{stacks-project:flat.tex:local-snapshot-20261011; GNU FDL 1.2 or later}
\begin{document}
\maketitle
\noindent\textit{Editorial note.} All delivered mathematical text is the author's own, copied from the pinned Stacks Project chapter file flat.tex; no statement or proof step was rewritten, reordered or removed.\par\medskip
\noindent\textit{Editorial note.} The preamble is the chapter's own retained minimal preamble (sources/stacks-flat/retained-preamble.tex), a verbatim excerpt of the native Xy-pic setup and of the native theorem declarations and mathematical macros, followed by the project own definitions of the reference, slogan and history display environments (verbatim package plus three newenvironment lines) that the retained file carries at its end and that the project itself typesets as comments. No macro was invented or redefined, and the wrapper class and the amsmath/amssymb/amsthm/graphicx base packages are supplied by the assembler.\par\medskip
\noindent\textit{Editorial note.} The counted statement and proof contain no cross reference, no citation and no included file, so the delivered body is self contained.\par\medskip
\noindent\textit{Editorial note.} Canonical location of the counted proposition: Stacks tag 0GTX (sources/stacks-flat/pinned-original-source.tex line 13535).\par\medskip
\noindent\textit{Allowed standard prerequisites (editorial).}\begin{itemize}
\item The derived twist L eta\_I and the complex eta\_I M\textasciicircum{}bullet. Cohomology, Lemma labelled cohomology-lemma-Leta, cited at proof lines 13546-13548. Identifies eta\_I M\textasciicircum{}bullet as a representative of L eta\_I M.
\item Comparison of a bounded complex of finite free modules over an affine with the given object. Derived Categories of Schemes, Lemma labelled perfect-lemma-affine-compare-bounded, cited at proof lines 13558-13562. Produces the bounded complex N\textasciicircum{}bullet of finite free A-modules representing M\textasciicircum{}bullet over the chart.
\item Principal ideals I\_i(N\textasciicircum{}bullet, f) and the behaviour of eta\_f. More on Algebra, Lemma labelled more-algebra-lemma-ideal-well-defined, cited at proof lines 13562-13564. Gives that eta\_f N\textasciicircum{}bullet is bounded with finite locally free terms.
\item Stalks of the twist. Same chapter, Lemma labelled lemma-eta-stalks, cited at proof lines 13564-13566. Identifies the quasi-coherent module Q\textasciicircum{}i|\_U with the one attached to eta\_f N\textasciicircum{}i.
\end{itemize}\medskip
\ifdefined\hypertarget\hypertarget{sourcepointer-1}{}\fi
\noindent\textit{Source reference.} Original reference 0F8S: sources/stacks-flat/pinned-original-source.tex, source lines 13460--13460. Same-chapter pointer into the author environment situation at source line 13460; invoked at proof line None of the counted proof. It is retained only as a navigation pointer; every load-bearing step is inside the delivered span.\par\medskip
\ifdefined\hypertarget\hypertarget{sourcepointer-2}{}\fi
\noindent\textit{Source reference.} Original lemma 0F8Q: sources/stacks-chow/tags.txt, source lines 19050--19050. Cross-chapter standard prerequisite located in the project tag list staged in the local snapshot; invoked at proof line 13547 of the counted proof. It is retained only as a navigation pointer; every load-bearing step is inside the delivered span.\par\medskip
\ifdefined\hypertarget\hypertarget{sourcepointer-3}{}\fi
\noindent\textit{Source reference.} Original lemma 06Z0: sources/stacks-chow/tags.txt, source lines 8904--8904. Cross-chapter standard prerequisite located in the project tag list staged in the local snapshot; invoked at proof line 13560 of the counted proof. It is retained only as a navigation pointer; every load-bearing step is inside the delivered span.\par\medskip
\ifdefined\hypertarget\hypertarget{sourcepointer-4}{}\fi
\noindent\textit{Source reference.} Original lemma 0GSS: sources/stacks-chow/tags.txt, source lines 20942--20942. Cross-chapter standard prerequisite located in the project tag list staged in the local snapshot; invoked at proof line 13562 of the counted proof. It is retained only as a navigation pointer; every load-bearing step is inside the delivered span.\par\medskip
\ifdefined\hypertarget\hypertarget{sourcepointer-5}{}\fi
\noindent\textit{Source reference.} Original lemma 0F8M: sources/stacks-flat/pinned-original-source.tex, source lines 13412--13412. Same-chapter pointer into the author environment lemma at source line 13412; invoked at proof line 13566 of the counted proof. It is retained only as a navigation pointer; every load-bearing step is inside the delivered span.\par\medskip
\begin{lemma}
\label{lemma-good-triple-bdd-loc-free}
Let $X, D, \mathcal{I}, M$ be as in
Situation \ref{situation-complex-and-divisor}.
Assume $(X, D, M)$ is a good triple.
If there exists a locally bounded complex $\mathcal{M}^\bullet$
of finite locally free $\mathcal{O}_X$-modules representing $M$,
then there exists a locally bounded complex $\mathcal{Q}^\bullet$
of finite locally free $\mathcal{O}_{X'}$-modules representing
$L\eta_\mathcal{I}M$.
\end{lemma}

\begin{proof}
By Cohomology, Lemma \ref{cohomology-lemma-Leta}
the complex $\mathcal{Q}^\bullet = \eta_\mathcal{I}\mathcal{M}^\bullet$
represents $L\eta_\mathcal{I}M$. To check that this complex
is locally bounded and consists of finite locally free, we may work
affine locally. Then the boundedness is clear.
Choose an affine chart $(U, A, f, M^\bullet)$ for $(X, D, M)$
such that the ideals $I_i(M^\bullet, f)$ are principal and such that
$\mathcal{M}^i|_U$ is finite free for each $i$. By our assumption
that $(X, D, M)$ is a good triple we can do this. Writing
$N^i = \Gamma(U, \mathcal{M}^i|_U)$ we get a bounded
complex $N^\bullet$ of finite free $A$-modules representing the
same object in $D(A)$ as the complex $M^\bullet$ (by
Derived Categories of Schemes, Lemma
\ref{perfect-lemma-affine-compare-bounded}).
Then $I_i(N^\bullet, f)$ is a principal ideal for all $i$
by More on Algebra, Lemma \ref{more-algebra-lemma-ideal-well-defined}.
Hence the complex $\eta_fN^\bullet$ is a bounded complex
of finite locally free $A$-modules. Since
$\mathcal{Q}^i|_U$ is the quasi-coherent $\mathcal{O}_U$-module
corresponding to $\eta_fN^i$ by Lemma \ref{lemma-eta-stalks} we conclude.
\end{proof}

\end{document}
